\documentclass[11pt,a4paper]{article}

\usepackage{amsmath,amssymb,amsthm,mathtools}
\usepackage{booktabs}
\usepackage{graphicx}
\usepackage[margin=2.7cm]{geometry}
\usepackage[colorlinks=true,linkcolor=blue,urlcolor=blue]{hyperref}
\graphicspath{{../figures/}}
\hypersetup{
  pdftitle={Static Inverse-Gamma Variance Mixtures, Student-t Option Valuation,
    Risk Premia, and Stationary Laws of Stochastic-Volatility Diffusions},
  pdfsubject={Student-t option valuation and physical-to-pricing risk premia}
}

\newtheorem{proposition}{Proposition}[section]
\newtheorem{remark}{Remark}[section]
\newtheorem{assumption}{Assumption}[section]

\DeclareMathOperator{\IG}{IG}
\DeclareMathOperator{\GammaDist}{Gamma}
\DeclareMathOperator{\Var}{Var}
\DeclareMathOperator{\Cov}{Cov}
\DeclareMathOperator{\E}{E}

\newcommand{\dd}{\,\mathrm{d}}
\newcommand{\ind}{\mathbf{1}}
\newcommand{\cdf}{G}
\newcommand{\cM}{\mathcal{M}}
\newcommand{\doi}[1]{\href{https://doi.org/#1}{doi:#1}}

\title{Static Inverse-Gamma Variance Mixtures, Student-$t$ Option Valuation,\\
Risk Premia, and Stationary Laws of Stochastic-Volatility Diffusions}
\author{}
\date{Draft: August 2026}

\begin{document}

\maketitle

\begin{abstract}
This note gives a self-contained formulation of a static inverse-gamma normal
variance mixture and the resulting Student-$t$ model for European option
valuation.  The terminal asset price is non-negative, its drift is fixed by the
market forward, and calls, puts, put--call parity, deltas, gammas, and strike
derivatives are obtained in closed form.  We then construct an explicit
physical-to-pricing measure change.  A power--precision kernel preserves the
inverse-gamma family, while a conditionally compensated Gaussian kernel
introduces asymmetric return risk without relying on the nonexistent Student
moment-generating function.  The three parameters have interpretable smile
effects: variance level, tail curvature, and skew.  The notation is chosen to
match the symmetric implementation in
\texttt{stochvolmodels.fitters.tdist}; the accompanying
\texttt{risk\_premium\_smile\_examples.py} implements the $p,\eta,q$
extension.  A separate section derives the stationary laws of the
multiplicative-variance, $3/2$, and Heston/CIR variance diffusions.  The
stationary laws motivate static return mixtures, but they are not identified
with the finite-horizon return laws of the continuous-time diffusions.
\end{abstract}

\tableofcontents

\section{Scope and conventions}

The note separates three related but different constructions.
\begin{enumerate}
    \item A \emph{static terminal model}: a conditional variance is drawn from
    an inverse-gamma law and a return shock is conditionally normal.  The
    resulting marginal shock is Student-$t$.  This construction is used for
    option valuation one maturity at a time.
    \item A \emph{physical-to-pricing measure change}: an admissible likelihood
    ratio changes the inverse-gamma shape and scale and can couple large
    variance to signed terminal returns.  This construction defines risk
    premia at one maturity and produces skewed option smiles.
    \item \emph{Continuous-time motivation}: positive variance diffusions have
    inverse-gamma or gamma stationary laws.  Their stationary
    marginal variance may motivate a one-step mixture, but the exact
    finite-horizon return mixes over integrated variance.
\end{enumerate}

Throughout, the inverse-gamma convention is
\begin{equation}
  V\sim\IG(\alpha,\beta),\qquad
  g_{\IG}(v)=\frac{\beta^\alpha}{\Gamma(\alpha)}
  v^{-\alpha-1}\exp\left(-\frac{\beta}{v}\right),\qquad v>0,
  \label{eq:ig_density}
\end{equation}
where $\alpha>0$ is the shape and $\beta>0$ is the scale.  When they exist,
\begin{equation}
  \E[V]=\frac{\beta}{\alpha-1},\qquad
  \Var[V]=\frac{\beta^2}{(\alpha-1)^2(\alpha-2)}.
  \label{eq:ig_moments}
\end{equation}
The gamma convention uses a rate $\lambda>0$:
\begin{equation}
  V\sim\GammaDist(\alpha,\lambda),\qquad
  g_\Gamma(v)=\frac{\lambda^\alpha}{\Gamma(\alpha)}
  v^{\alpha-1}e^{-\lambda v},\qquad v>0.
  \label{eq:gamma_density}
\end{equation}

For option valuation at maturity $T$, let $D_T$ be the market discount factor,
$F_T$ the market forward, and
\begin{equation}
  \bar S_0(T):=D_TF_T
  \label{eq:prepaid_forward}
\end{equation}
the prepaid forward.  This convention incorporates deterministic carry and
dividends in $F_T$.  The chain-pricing implementation passes
\texttt{forward * discfactor} as its argument named \texttt{spot}; consequently
that argument is precisely $\bar S_0(T)$ in the formulation below.

\section{Related literature and positioning}
\label{sec:related_literature}

The main ingredients of the model have substantial precedents, although they
are usually assembled in different ways.

\begin{description}
  \item[Normal--inverse-gamma foundation.]
  Praetz \cite{Praetz1972} derived Student-$t$ laws for share-price changes by
  mixing conditional normal returns over an inverse-gamma variance.  This is
  the distributional foundation of the static construction below, but that
  work does not address risk-neutral option valuation or the lower price
  boundary.

  \item[Arithmetic Student returns.]
  Pinn \cite{Pinn2000} gives the closest early option model based on additive
  Student-$t$ price changes and derives European-call costs and
  minimum-variance hedge functions.  Its pricing principle is
  variance-optimal hedging in an incomplete market; negative underlying values
  remain possible and there is no censoring-induced atom at zero.  Vince
  \cite{Vince1992} separately observed that imposing a zero price bound changes
  the terminal mean and can disturb put--call parity, and proposed a numerical
  compensating shift.  That construction is based on log-Student moves and
  does not solve the censored martingale equation used here.

  \item[Log-Student valuation and Greeks.]
  Cassidy, Hamp, and Ouyed \cite{CassidyHampOuyed2010} construct martingale-
  normalized log-Student option formulas with a cap or truncated support.
  Such upper-tail control is necessary because a Student variable has no
  finite positive exponential moment.  Their follow-up paper
  \cite{CassidyHampOuyed2013} adds option Greeks and lower-floor extensions,
  including the sensitivity of the martingale normalizer.  In contrast,
  \eqref{eq:terminal_asset} is arithmetic: for $\nu>1$ it needs no upper cap,
  and its entire censored lower tail becomes an atom at $S_T=0$.

  \item[Inverse-gamma activity and geometric prices.]
  Castelli, Leonenko, and Shchestyuk
  \cite{CastelliLeonenkoShchestyuk2017} use a stationary reciprocal-gamma
  activity process to generate Student-like returns and an option-pricing
  formula.  Leonenko, Liu, and Shchestyuk
  \cite{LeonenkoLiuShchestyuk2025} develop European prices, analytic Delta,
  Gamma, Vega, and Theta, and delta-hedging experiments for Student
  fractal-activity-time geometric Brownian motion.  Conditional
  Black--Scholes quantities are averaged over activity time.  Ivanov
  \cite{Ivanov2023} obtains related option formulas when generalized
  Black--Scholes--Merton coefficients depend on gamma or inverse-gamma random
  variables.  These exponential constructions keep the asset strictly
  positive and have neither a zero atom nor the endogenous censored location
  $A_T$.

  \item[Generalized-hyperbolic skew Student law.]
  The normal mean--variance mixture in \eqref{eq:q_joint_law}, and hence the
  marginal density \eqref{eq:skew_t_density}, is a parameterization of the
  established generalized-hyperbolic skew Student-$t$ distribution studied by
  Aas and Hob\ae k Haff \cite{AasHaff2006}.  They emphasize its polynomial
  tail on one side and exponential tail on the other.  The known marginal
  family is not the proposed contribution here; the added structure is the
  explicit physical-to-pricing kernel, the floored martingale location,
  option valuation, and risk-premium Greeks.

  \item[Pricing kernels and measure selection.]
  The power--precision likelihood ratio developed below belongs to the broad
  class of positive pricing kernels that preserve a chosen parametric family.
  In an incomplete market it selects one equivalent martingale measure among
  many; it is not forced by no-arbitrage alone.  Minimal Hellinger and more
  general minimal $f^q$ martingale measures provide principled selection
  criteria in related semimartingale settings
  \cite{ChoulliStrickerLi2007,JeanblancKloppelMiyahara2007}.  For dynamic
  diffusions, a formally specified drift change must additionally produce a
  true density process and preserve the relevant boundary behavior; the
  one-dimensional criteria in Desmettre, Leobacher, and Rogers
  \cite{DesmettreLeobacherRogers2021} make explicit why a local Girsanov
  calculation alone is insufficient.
  Here the superscript in $f^q$ is standard notation from that literature and
  is unrelated to the return-skew parameter $q$ in
  \eqref{eq:conditional_return_kernel}.

  \item[Stationary SV laws.]
  Inverse-gamma stationarity for the multiplicative/GARCH diffusion goes back
  to Nelson \cite{Nelson1990}.  The $3/2$ diffusion is represented in the
  nonlinear model of Ahn and Gao \cite{AhnGao1999}; its reciprocal is a
  square-root diffusion.  The CIR and Heston square-root specifications instead
  have gamma stationary laws \cite{CIR1985,Heston1993}.  Dynamic GARCH-diffusion
  option valuation has been developed through moments of \emph{integrated}
  variance \cite{BaroneAdesiRasmussenRavanelli2005}, not by substituting the
  stationary instantaneous variance into a one-period Student mixture.
\end{description}

Thus Student-$t$ option valuation, inverse-gamma mixing, martingale
normalization, lower boundaries, and Greeks are not individually new.  In the
targeted literature search, however, no exact published match was located for
their joint use in
\begin{equation}
  S_T=\bar S_0(T)(A_T+X_T)^+,
  \qquad X_T\ \hbox{centered Student-}t,
  \label{eq:literature_position_model}
\end{equation}
with $A_T$ determined \emph{after} censoring from the martingale condition,
uncapped arithmetic upside, a genuine zero-price atom, exact calls and puts,
and Greeks that include the implicit response of $A_T$.  The partial-moment
option identity itself is standard; the potentially distinctive contribution
is this complete, internally consistent package.  This conclusion records the
outcome of a targeted literature review and is not a formal claim of priority.

\section{Static inverse-gamma variance mixtures}

\subsection{The normal variance mixture}

\begin{assumption}[Static conditional return]
For a fixed horizon with scale factor $c>0$,
\begin{equation}
  X=\sqrt{cV}\,Z,\qquad Z\sim N(0,1),\qquad
  Z\perp V,\qquad V\sim\IG(\alpha,\beta),
  \label{eq:normal_ig_mixture}
\end{equation}
Equivalently, $X\mid V\sim N(0,cV)$, where the second normal parameter is
the variance.
\end{assumption}

\begin{proposition}[Student-$t$ marginal]
Under \eqref{eq:normal_ig_mixture},
\begin{equation}
  f_X(x)=
  \frac{\Gamma(\alpha+\tfrac12)}
       {\Gamma(\alpha)\sqrt{2\pi c\beta}}
  \left(1+\frac{x^2}{2c\beta}\right)^{-(\alpha+1/2)},
  \label{eq:mixture_density}
\end{equation}
or, equivalently,
\begin{equation}
  X\overset{d}{=}sZ_\nu,\qquad
  Z_\nu\sim t_\nu,\qquad
  \nu=2\alpha,\qquad
  s^2=\frac{c\beta}{\alpha}.
  \label{eq:student_mapping}
\end{equation}
\end{proposition}

\begin{proof}
Integrating the conditional normal density against \eqref{eq:ig_density} gives
\begin{align*}
f_X(x)
&=\frac{\beta^\alpha}{\Gamma(\alpha)\sqrt{2\pi c}}
  \int_0^\infty
  v^{-\alpha-3/2}
  \exp\left[-\frac{\beta+x^2/(2c)}{v}\right]\dd v\\
&=\frac{\beta^\alpha\Gamma(\alpha+\tfrac12)}
        {\Gamma(\alpha)\sqrt{2\pi c}}
  \left(\beta+\frac{x^2}{2c}\right)^{-(\alpha+1/2)},
\end{align*}
which is \eqref{eq:mixture_density}.  Substitution of $\nu=2\alpha$ and
$s^2=c\beta/\alpha$ yields the standard location-scale Student density.
\end{proof}

If $\alpha>1$, then
\begin{equation}
  \Var[X]=c\E[V]=\frac{c\beta}{\alpha-1}
  =s^2\frac{\nu}{\nu-2}.
  \label{eq:variance_match_ig}
\end{equation}
Thus, for an annualized volatility parameter $\sigma$ and $c=T$, the
implementation convention
\begin{equation}
  \Var[X_T]=\sigma^2T
  \label{eq:code_variance_target}
\end{equation}
is obtained with
\begin{equation}
  s_T=\sigma\sqrt{T\frac{\nu-2}{\nu}},\qquad \nu>2.
  \label{eq:code_student_scale}
\end{equation}
This is the quantity returned by \texttt{compute\_upsilon}; despite its name,
it is a Student scale, not a precision.

\subsection{Density, distribution, and partial first moment}

Set
\begin{equation}
  c_\nu:=\frac{\Gamma((\nu+1)/2)}
  {\sqrt{\pi\nu}\,\Gamma(\nu/2)}.
\end{equation}
For the centered shock $X_T=s_TZ_\nu$,
\begin{equation}
  f_T(x)=\frac{c_\nu}{s_T}
  \left(1+\frac{x^2}{\nu s_T^2}\right)^{-(\nu+1)/2}.
  \label{eq:student_density}
\end{equation}
Let
\begin{equation}
  \cdf_T(x):=\Pr(X_T\le x),\qquad
  H_T(x):=\int_{-\infty}^x u f_T(u)\dd u,
  \label{eq:cdf_lower_moment}
\end{equation}
and define the upper partial first moment by
\begin{equation}
  \cM_T(x):=\int_x^\infty u f_T(u)\dd u=-H_T(x).
  \label{eq:upper_moment}
\end{equation}
For $\nu>1$,
\begin{equation}
  H_T(x)=
  -s_Tc_\nu\frac{\nu}{\nu-1}
  \left(1+\frac{x^2}{\nu s_T^2}\right)^{-(\nu-1)/2}.
  \label{eq:student_lower_moment}
\end{equation}
The implementation functions \texttt{cdf\_tdist} and
\texttt{cum\_mean\_tdist} compute $\cdf_T$ and $H_T$, respectively, with
their location argument set to zero in the terminal pricer.

For a generic moment order $r$, the condition is
\begin{equation}
  \E[|X_T|^r]<\infty\quad\Longleftrightarrow\quad r<\nu.
  \label{eq:student_moment_condition}
\end{equation}
Linear option payoffs require $\nu>1$, variance matching requires $\nu>2$,
and finite kurtosis requires $\nu>4$.  The implementation deliberately imposes
$\nu>2$ because \eqref{eq:code_student_scale} uses a finite variance.

\subsection{Transforms}

The moment-generating function of a Student variable satisfies
$\E[e^{tX}]=\infty$ for every real $t\ne0$.  The relevant transform is the characteristic
function, obtained from the Laplace transform of the mixing variance:
\begin{equation}
  \phi_X(u)=\E[e^{iuX}]=\E[e^{-cu^2V/2}]
  =L_V(cu^2/2).
\end{equation}
For \eqref{eq:ig_density},
\begin{equation}
  L_V(\ell)=\E[e^{-\ell V}]
  =\frac{2(\beta \ell)^{\alpha/2}}{\Gamma(\alpha)}
  K_\alpha(2\sqrt{\beta \ell}),\qquad \ell>0,
  \label{eq:ig_laplace}
\end{equation}
where $K_\alpha$ is the modified Bessel function of the second kind.

\section{Physical-to-pricing measure changes and risk premia}
\label{sec:risk_premia}

\subsection{One-period pricing densities}

Let the physical law at a fixed maturity be
\begin{equation}
  V\sim\IG(\alpha_{\mathbb P},\beta_{\mathbb P}),\qquad
  X\mid V,\mathbb P\sim N(0,cV).
  \label{eq:physical_joint_law}
\end{equation}
A one-period change of measure is a strictly positive random variable
\begin{equation}
  \Lambda_T=\frac{\dd\mathbb Q}{\dd\mathbb P}\bigg|_{\mathcal F_T},
  \qquad
  \E_{\mathbb P}[\Lambda_T]=1.
  \label{eq:terminal_rn_density}
\end{equation}
For a terminal price $S_T=B_T(A_T+X_T)^+$ with an arbitrary positive price
scale $B_T$, the pricing law must also satisfy
\begin{equation}
  \E_{\mathbb P}[\Lambda_TS_T]
  =\E_{\mathbb Q}[S_T]=F_T.
  \label{eq:pricing_measure_martingale}
\end{equation}
The second restriction fixes $A_T$ after the distributional parameters have
been selected.  It does not uniquely fix $\Lambda_T$: the one-period market is
incomplete.

\subsection{Why the marginal Student Esscher transform fails}

The conventional Esscher candidate
\begin{equation}
  \Lambda_T^{\rm Esscher}
  =\frac{e^{tX_T}}{\E_{\mathbb P}[e^{tX_T}]}
  \label{eq:invalid_student_esscher}
\end{equation}
does not exist for $t\ne0$, because the Student moment-generating function is
infinite on both non-zero half-axes.  This rules out an exponential tilt of
the \emph{marginal} Student shock.  It does not rule out every exponential
factor on the enlarged state space $(X,V)$: conditional Gaussian
normalization makes the construction below finite.

\subsection{Inverse-gamma-preserving power--precision kernel}

For parameters $p$ and $\eta$, define
\begin{equation}
  \Lambda_{p,\eta}^{V}
  :=
  \frac{V^p e^{-\eta/V}}
       {\E_{\mathbb P}[V^p e^{-\eta/V}]}.
  \label{eq:variance_pricing_kernel}
\end{equation}
The normalizing moment is
\begin{equation}
  \E_{\mathbb P}[V^p e^{-\eta/V}]
  =
  \frac{\beta_{\mathbb P}^{\alpha_{\mathbb P}}
        \Gamma(\alpha_{\mathbb P}-p)}
       {\Gamma(\alpha_{\mathbb P})
        (\beta_{\mathbb P}+\eta)^{\alpha_{\mathbb P}-p}},
  \label{eq:variance_kernel_normalizer}
\end{equation}
which is finite precisely when
\begin{equation}
  p<\alpha_{\mathbb P},\qquad
  \eta>-\beta_{\mathbb P}.
  \label{eq:static_kernel_admissibility}
\end{equation}
Multiplying \eqref{eq:ig_density} by \eqref{eq:variance_pricing_kernel}
preserves the inverse-gamma family:
\begin{equation}
  V\mid\mathbb Q\sim\IG(a,b),\qquad
  a:=\alpha_{\mathbb P}-p,\qquad
  b:=\beta_{\mathbb P}+\eta.
  \label{eq:q_ig_parameters}
\end{equation}
Thus, in the symmetric $q=0$ subfamily, $p$ changes the Student degrees of
freedom $\nu_{\mathbb Q}=2a$, while $\eta$ changes the variance scale.  When
$q\ne0$, the polynomial-side moment index is instead $a$, as made explicit
in \eqref{eq:skew_moment_condition}.  In particular,
\begin{equation}
  \bar v_{\mathbb Q}:=\E_{\mathbb Q}[V]
  =\frac{b}{a-1},\qquad a>1.
  \label{eq:q_mean_variance}
\end{equation}

\subsection{Conditionally compensated return kernel}

Define a second likelihood ratio conditionally on $V$:
\begin{equation}
  \Lambda_q^{X\mid V}
  :=
  \exp\left(\frac{qX}{c}-\frac{q^2V}{2c}\right).
  \label{eq:conditional_return_kernel}
\end{equation}
Since $X\mid V,\mathbb P\sim N(0,cV)$,
\begin{equation}
  \E_{\mathbb P}[\Lambda_q^{X\mid V}\mid V]=1.
\end{equation}
Consequently, the product
\begin{equation}
  \Lambda_{p,\eta,q}
  =\Lambda_{p,\eta}^{V}\Lambda_q^{X\mid V}
  \label{eq:combined_pricing_kernel}
\end{equation}
is a valid, strictly positive one-period density under
\eqref{eq:static_kernel_admissibility}.  Under the resulting pricing law,
\begin{equation}
  V\sim\IG(a,b),\qquad
  X\mid V,\mathbb Q\sim N(qV,cV),
  \qquad
  Z_{\mathbb Q}:=\frac{X-qV}{\sqrt{cV}}\sim N(0,1),
  \quad Z_{\mathbb Q}\perp_{\mathbb Q}V.
  \label{eq:q_joint_law}
\end{equation}
The sign convention is deliberate: $q<0$ associates large-variance states
with negative returns and generates an equity-like left skew.

For $q\ne0$, direct integration gives the skewed density
\begin{align}
  f_{\mathbb Q}(x)
  &=
  \frac{2b^a}{\Gamma(a)\sqrt{2\pi c}}\,
  e^{qx/c}
  \left(\frac{q^2}{x^2+2bc}\right)^{(a+1/2)/2}
  \nonumber\\
  &\hspace{2cm}\times
  K_{a+1/2}\left(
    \frac{|q|}{c}\sqrt{x^2+2bc}
  \right).
  \label{eq:skew_t_density}
\end{align}
Its continuous $q\to0$ limit is the symmetric Student density
\eqref{eq:mixture_density} with $\alpha=a$ and $\beta=b$.  The construction is therefore an
inverse-gamma normal mean--variance mixture, not a marginal Student Esscher
transform.

The moment implications are important.  For $q=0$,
$\E_{\mathbb Q}[|X|^r]<\infty$ if and only if $r<2a$.  For $q\ne0$, the
$qV$ term dominates one heavy tail, so
\begin{equation}
  \E_{\mathbb Q}[|X|^r]<\infty
  \quad\Longleftrightarrow\quad r<a.
  \label{eq:skew_moment_condition}
\end{equation}
The domain $a>1$ is a uniform sufficient condition for a finite first
absolute moment and linear valuation for every sign of $q$; $a>2$ is a
uniform sufficient condition for finite variance.  The exact symmetric
$q=0$ thresholds are weaker: a finite positive-part value requires
$a>1/2$, and finite variance requires $a>1$.  For $q>0$, linear valuation
requires $a>1$; for $q<0$, the payoff-relevant right tail is exponentially
tempered even when the first absolute moment is infinite.  When $a>2$,
\begin{equation}
  \Var_{\mathbb Q}[X]
  =c\bar v_{\mathbb Q}
   +q^2\Var_{\mathbb Q}[V]
  =c\bar v_{\mathbb Q}
   +\frac{q^2\bar v_{\mathbb Q}^2}{a-2}.
  \label{eq:q_total_variance}
\end{equation}

\subsection{Valuation by conditional Bachelier integration}

Introduce the normal positive-part function
\begin{equation}
  \mathcal L(m,s):=m\Phi(m/s)+s\phi(m/s),\qquad s>0.
  \label{eq:normal_positive_part}
\end{equation}
Conditional on $V$, it equals
$\E[(m+sZ)^+]$.  The martingale location is therefore the unique solution
of
\begin{equation}
  \E_{\mathbb Q}\!\left[
    \mathcal L(A_T+qV,\sqrt{cV})
  \right]
  =\frac{F_T}{B_T}.
  \label{eq:skew_martingale_location}
\end{equation}
For strike $K$, define
\begin{equation}
  m_K(V):=A_T-\frac{K}{B_T}+qV,\qquad
  s(V):=\sqrt{cV}.
\end{equation}
The call and put prices are
\begin{align}
  C_0(K)
  &=D_TB_T\,\E_{\mathbb Q}
    [\mathcal L(m_K(V),s(V))],
  \label{eq:skew_call_price}
  \\
  P_0(K)
  &=C_0(K)-D_T(F_T-K).
  \label{eq:skew_put_price}
\end{align}
The put expression includes the zero-state payoff through parity.  Direct
conditional integration gives the same result when the lower tail is handled
with the sign-reversed normal mean.

The one-dimensional expectation is stable under the precision change
\begin{equation}
  U:=\frac{b}{V}\sim\GammaDist(a,1),
  \label{eq:gamma_precision_quadrature}
\end{equation}
so generalized Gauss--Laguerre quadrature applies.  At $q=0$,
\eqref{eq:skew_call_price} agrees with the closed-form Student partial-moment
formula below.  This equality is a useful independent implementation check.

\subsection{Strike derivatives and risk-premium Greeks}

Write
\begin{equation}
  z_K(V):=\frac{m_K(V)}{s(V)},\qquad
  z_0(V):=\frac{A_T+qV}{s(V)}.
\end{equation}
Differentiating under the integral gives
\begin{align}
  \frac{\partial C_0}{\partial K}
  &=-D_T\,\E_{\mathbb Q}[\Phi(z_K(V))],
  \label{eq:skew_call_strike_delta}
  \\
  \frac{\partial^2 C_0}{\partial K^2}
  &=\frac{D_T}{B_T}\,
    \E_{\mathbb Q}\left[
      \frac{\phi(z_K(V))}{s(V)}
    \right].
  \label{eq:skew_call_strike_gamma}
\end{align}
Hence the call is decreasing and convex in strike, and the second derivative
recovers the continuous terminal density away from zero.

Risk-premium Greeks must include the response of $A_T$.  The following
formulas use the uniform domain $a>1$, which ensures the required first
moments for $q=0$ and $q>0$.  For
$V\sim\IG(a,b)$, the centered likelihood scores are
\begin{equation}
  s_p(V)=\log(V/b)+\psi(a),\qquad
  s_\eta(V)=\frac{a}{b}-\frac{1}{V},
  \label{eq:ig_risk_premium_scores}
\end{equation}
where $\psi$ is the digamma function.  For
$\vartheta\in\{p,\eta\}$,
\begin{equation}
  A_\vartheta
  =
  -\frac{
    \E_{\mathbb Q}[
      \mathcal L(A_T+qV,s(V))s_\vartheta(V)
    ]}
    {\E_{\mathbb Q}[\Phi(z_0(V))]}.
  \label{eq:location_p_eta_greek}
\end{equation}
The corresponding price Greek is
\begin{equation}
  \frac{\partial C_0}{\partial\vartheta}
  =
  D_TB_T\left\{
    \E_{\mathbb Q}[
      \mathcal L(m_K(V),s(V))s_\vartheta(V)
    ]
    +A_\vartheta\E_{\mathbb Q}[\Phi(z_K(V))]
  \right\}.
  \label{eq:call_p_eta_greek}
\end{equation}
For the skew parameter,
\begin{align}
  A_q
  &=
  -\frac{\E_{\mathbb Q}[V\Phi(z_0(V))]}
         {\E_{\mathbb Q}[\Phi(z_0(V))]},
  \label{eq:location_q_greek}
  \\
  \frac{\partial C_0}{\partial q}
  &=
  D_TB_T\,\E_{\mathbb Q}[
    (A_q+V)\Phi(z_K(V))
  ].
  \label{eq:call_q_greek}
\end{align}
Equations \eqref{eq:location_p_eta_greek}--\eqref{eq:call_q_greek} hold with
$B_T,F_T,D_T,c$ fixed.  They show why bumping a risk-premium parameter while
holding $A_T$ fixed is inconsistent with the forward.

\subsection{Power and polynomial alternatives}

Power kernels are therefore possible and particularly convenient on the
latent variance.  The factor
\[
  V^p e^{-\eta/V}
\]
is positive and preserves the inverse-gamma family.  A direct asset-power
kernel is less satisfactory in
the floored model.  If it is proportional to $S_T^\gamma$, then
$\gamma>0$ assigns zero density to the positive-probability state $S_T=0$
and destroys equivalence, while $\gamma<0$ is singular at that atom.  A
strictly positive shift such as $(\epsilon+S_T)^\gamma$ avoids this defect
but loses the family-preserving formulas.

Polynomial kernels are possible only after positivity and integrability are
enforced.  An odd polynomial such as
$1+a_1X+a_3X^3$ changes sign in a tail and cannot be a
Radon--Nikodym density.  A strictly positive even polynomial can be
normalized when all required Student moments exist; for a degree-$2m$
polynomial in the symmetric model this requires $2m<2\alpha_{\mathbb P}$.
It generally produces a mixture of components rather than one
inverse-gamma law, is sensitive to high moments, and gives less transparent
smile controls.  A kernel such as
$\epsilon+h(X)^2$ with $\epsilon>0$ is strictly positive, and exponentials
of bounded functions or bounded transforms also preserve positivity.
These alternatives similarly sacrifice the closed family.  For these reasons,
\eqref{eq:combined_pricing_kernel} is the recommended parsimonious
formulation.

\subsection{Comparative statics of \texorpdfstring{$p,\eta,q$}{p, eta, q}}
\label{sec:risk_premium_comparative_statics}

The numerical illustration uses
\begin{equation}
  F_T=D_T=T=c=B_T=1,\qquad
  \alpha_{\mathbb P}=4,\qquad
  \beta_{\mathbb P}=0.12,\qquad
  \E_{\mathbb P}[V]=0.04.
  \label{eq:risk_premium_baseline}
\end{equation}
For every curve, $A_T$ is re-solved from
\eqref{eq:skew_martingale_location}.  The baseline
$(p,\eta,q)=(0,0,0)$ has an unconditional shock standard deviation of
$20\%$.  Its Black implied-volatility curve is nevertheless not flat:
arithmetic returns, log-moneyness quotation, heavy tails, and the lower
floor already produce a modest negative Black-volatility skew.

\begin{figure}[htbp]
  \centering
  \includegraphics[width=\textwidth]{risk_premium_smiles.pdf}
  \caption{One-year Black implied-volatility smiles for the three
  risk-premium parameters.  The left panel varies $p$ with $\eta=0$, so tail
  thickness and mean variance move together.  The middle panel varies
  $\eta$ at fixed $p=0$ and predominantly shifts the variance level.  The
  right panel varies $q$: negative $q$ couples high variance to negative
  returns and steepens the left skew.  The martingale location is re-solved
  for every curve.}
  \label{fig:risk_premium_smiles}
\end{figure}

A raw change in $p$ changes both tail thickness and mean variance:
\begin{equation}
  \nu_{\mathbb Q}=2(\alpha_{\mathbb P}-p),\qquad
  \E_{\mathbb Q}[V]
  =\frac{\beta_{\mathbb P}+\eta}
         {\alpha_{\mathbb P}-p-1}.
  \label{eq:raw_p_effect}
\end{equation}
To isolate tail shape, hold
$\bar v_{\mathbb Q}=\E_{\mathbb Q}[V]$ fixed and set
\begin{equation}
  \eta(p)
  =\bar v_{\mathbb Q}(\alpha_{\mathbb P}-p-1)
   -\beta_{\mathbb P}.
  \label{eq:variance_neutral_eta}
\end{equation}
At the physical baseline this simplifies to
$\eta(p)=-\bar v_{\mathbb P}p=-0.04p$.

\begin{figure}[htbp]
  \centering
  \includegraphics[width=0.82\textwidth]
    {p_tail_premium_fixed_variance.pdf}
  \caption{The $p$ premium with $\eta=-0.04p$, which holds
  $\E_{\mathbb Q}[V]=4\%$ fixed.  The ATM level moves little, while $p>0$
  fattens the tails and raises wing curvature.  This variance-neutral
  parameterization separates the tail premium from the scale premium.}
  \label{fig:p_fixed_variance}
\end{figure}

For quantitative comparison let $k=\log(K/F_T)$ and define
\begin{equation}
  \operatorname{Slope}_{\rm ATM}
  :=
  \frac{\sigma_{\rm IV}(0.025)-\sigma_{\rm IV}(-0.025)}{0.05}.
  \label{eq:atm_slope_definition}
\end{equation}
The wing diagnostics are
\begin{align}
  \operatorname{RR}^{(k)}_{0.25}
  &:=
  \sigma_{\rm IV}(0.25)-\sigma_{\rm IV}(-0.25),
  \\
  \operatorname{BF}^{(k)}_{0.25}
  &:=
  \tfrac12\{
    \sigma_{\rm IV}(-0.25)+\sigma_{\rm IV}(0.25)
  \}-\sigma_{\rm IV}(0).
  \label{eq:log_moneyness_rr_bf}
\end{align}
These are fixed-log-moneyness diagnostics, not $25$-delta market quotes.

\begin{table}[htbp]
\centering
\small
\begin{tabular}{@{}lrrrr@{}}
\toprule
Scenario
& ATM IV
& ATM slope
& $\operatorname{RR}^{(k)}_{0.25}$
& $\operatorname{BF}^{(k)}_{0.25}$ \\
\midrule
Baseline
& 0.1921 & -0.0961 & -0.0456 & 0.0123 \\
$p=-0.75$, $\eta=0$
& 0.1732 & -0.0867 & -0.0412 & 0.0112 \\
$p=+0.75$, $\eta=0$
& 0.2188 & -0.1089 & -0.0517 & 0.0138 \\
$p=-0.75$, $\eta=+0.030$, fixed $\E_{\mathbb Q}[V]$
& 0.1937 & -0.0970 & -0.0465 & 0.0103 \\
$p=+0.75$, $\eta=-0.030$, fixed $\E_{\mathbb Q}[V]$
& 0.1895 & -0.0945 & -0.0442 & 0.0153 \\
$\eta=-0.024$
& 0.1718 & -0.0859 & -0.0404 & 0.0133 \\
$\eta=+0.024$
& 0.2105 & -0.1052 & -0.0503 & 0.0115 \\
$q=-2$
& 0.1974 & -0.1612 & -0.0764 & 0.0122 \\
$q=+2$
& 0.1977 & -0.0354 & -0.0169 & 0.0120 \\
\bottomrule
\end{tabular}
\caption{Deterministic-quadrature comparative statics.  All volatilities are
in decimal units and $A_T$ is re-solved in every scenario.}
\label{tab:risk_premium_comparative_statics}
\end{table}

The intended division of labor is visible in
Figures~\ref{fig:risk_premium_smiles}--\ref{fig:p_fixed_variance} and
Table~\ref{tab:risk_premium_comparative_statics}.  The raw $p$ change moves
both ATM volatility and curvature.  The variance-neutral $p$ change
predominantly moves wing curvature.  The parameter $\eta$ predominantly
moves the ATM level.  The parameter $q$ predominantly moves slope:
changing $q$ from $-2$ to $+2$ changes the ATM slope from about $-0.161$ to
$-0.035$ while leaving ATM volatility near $19.8\%$.

The martingale compensation offsets the conditional mean without removing
asymmetry.  In the $q=-2$ and $q=+2$ scenarios,
\begin{equation}
  A_T=1.07974\quad\hbox{and}\quad A_T=0.92000,
\end{equation}
respectively.  The corresponding zero-state probabilities are approximately
$9.9\times10^{-4}$ and $3.1\times10^{-5}$, compared with
$2.1\times10^{-4}$ at the baseline.

A stable calibration parameterization replaces $\eta$ by
$\bar v_{\mathbb Q}$ and uses
\begin{equation}
  (\bar v_{\mathbb Q},p,q)
  =
  (\hbox{variance level},\hbox{tail shape},\hbox{asymmetry}).
\end{equation}
One can initialize $\bar v_{\mathbb Q}$ from ATM volatility, $q$ from a
risk reversal or local slope, and variance-neutral $p$ from a butterfly or
wing curvature; a joint refinement must re-solve $A_T$ at each evaluation.

\subsection{Dynamic interpretation and measure selection}

The terminal density \eqref{eq:combined_pricing_kernel} is valid for one
maturity but does not by itself define a time-consistent density process.
For the multiplicative diffusion \eqref{eq:multiplicative_variance}, use the
standard convention
\begin{equation}
  \frac{\dd\mathbb Q}{\dd\mathbb P}\bigg|_{\mathcal F_t}
  =
  \mathcal E\left(-\int_0^t\lambda_m(V_s)\dd W_s^{\mathbb P}\right),
  \qquad
  \lambda_m(v)
  =-\frac{\varepsilon_m}{2}
    \left(p+\frac{\eta}{v}\right).
  \label{eq:multiplicative_girsanov_kernel}
\end{equation}
Whenever this stochastic exponential is a true martingale and the boundary
conditions remain admissible, the drift under $\mathbb Q$ stays in the same
family with
\begin{equation}
  \kappa_m^{\mathbb Q}
  =\kappa_m^{\mathbb P}-\frac{\varepsilon_m^2p}{2},
  \qquad
  \kappa_m^{\mathbb Q}(\theta_m^{\mathbb Q})^2
  =
  \kappa_m^{\mathbb P}(\theta_m^{\mathbb P})^2
  +\frac{\varepsilon_m^2\eta}{2}.
  \label{eq:multiplicative_q_mapping}
\end{equation}
The stationary shape therefore changes by $-p$ and the stationary scale by
$+\eta$.  Positivity of mean reversion requires the stronger restriction
$p<\alpha_{\mathbb P}-1$.

For the $3/2$ diffusion \eqref{eq:three_halves_variance}, the corresponding
family-preserving market price is
\begin{equation}
  \lambda_3(v)
  =-\frac{\varepsilon_3}{2}
    \left(p\sqrt v+\frac{\eta}{\sqrt v}\right),
  \label{eq:three_halves_girsanov_kernel}
\end{equation}
with the same parameter mapping as
\eqref{eq:multiplicative_q_mapping} after replacing the subscript $m$ by
$3$.  Positive mean reversion now requires
$p<\alpha_{\mathbb P}-2$.  These inequalities are necessary, not sufficient:
true-martingale and boundary tests must be checked for the chosen initial law
and horizon \cite{DesmettreLeobacherRogers2021}.

The one-period parameter $q$ is likewise not automatically a
continuous-time equity-risk-premium coefficient.  Its dynamic realization
depends on the stock and variance Brownian drivers, their correlation, and
the remaining market price of risk in the incomplete stochastic-volatility
model.  Among the admissible measures, option calibration, an equilibrium
stochastic discount factor, minimal Hellinger distance, or another stated
criterion must select the pricing law
\cite{ChoulliStrickerLi2007,JeanblancKloppelMiyahara2007}.

\section{Arbitrage-consistent symmetric Student-\texorpdfstring{$t$}{t}
terminal model}

This section is the pricing-measure special case $q=0$ of
\eqref{eq:q_joint_law}.  Unless a measure is displayed explicitly, all
expectations in this section and the following Greek formulas are under
$\mathbb Q$.

\subsection{Non-negative terminal price and the martingale condition}

For one fixed maturity, specify
\begin{equation}
  S_T=\bar S_0(T)\,(A_T+X_T)^+,
  \qquad A_T=1+\mu_TT,
  \label{eq:terminal_asset}
\end{equation}
where $X_T$ is centered with density \eqref{eq:student_density}.  The default
boundary is
\begin{equation}
  x_T^*:=-A_T,
  \label{eq:default_boundary}
\end{equation}
and the model has an atom at zero of size
\begin{equation}
  p_T^{(0)}=\cdf_T(x_T^*).
  \label{eq:default_probability}
\end{equation}
This is a terminal zero-price mass, called a default probability in the code;
the static model does not define a first-passage time or a pathwise default
event.

The normalization in \eqref{eq:terminal_asset} is the one used by the code.
Consequently, $\sigma$ is the annualized standard deviation of an additive
arithmetic shock normalized by the prepaid forward $\bar S_0=D_TF_T$.  It is
neither a log-return volatility nor numerically identical to a Black--Scholes
volatility.  An equivalent forward-normalized representation is
\begin{equation}
  S_T=F_T(\widehat A_T+\widehat X_T)^+,
  \qquad
  \widehat A_T=D_TA_T,
  \qquad
  \widehat X_T=D_TX_T,
  \qquad
  \widehat\sigma_T=D_T\sigma_T.
  \label{eq:forward_normalization}
\end{equation}
It satisfies $\E[(\widehat A_T+\widehat X_T)^+]=1$.  Switching between the
two normalizations without rescaling $\sigma$ changes the model.

Risk neutrality requires $\E[S_T]=F_T$.  Since
\begin{equation}
  \E[(A_T+X_T)^+]
  =A_T[1-\cdf_T(-A_T)]-H_T(-A_T),
\end{equation}
and $F_T/\bar S_0(T)=D_T^{-1}$, the drift compensation is the unique solution
of
\begin{equation}
  A_T[1-\cdf_T(-A_T)]-H_T(-A_T)=D_T^{-1}.
  \label{eq:martingale_A}
\end{equation}
Equivalently, with $A_T=1+\mu_TT$,
\begin{equation}
  \mu_TT[1-\cdf_T(x_T^*)]
  -\cdf_T(x_T^*)-H_T(x_T^*)
  =D_T^{-1}-1.
  \label{eq:martingale_mu}
\end{equation}
Equation \eqref{eq:martingale_mu} is the root equation implemented by
\texttt{imply\_drift\_tdist}.  The left-hand side of
\eqref{eq:martingale_A} is continuous and strictly increasing in $A_T$, with
derivative $1-\cdf_T(-A_T)>0$, so the solution is unique.

\begin{remark}[Static, maturity-specific law]
The parameters $(\sigma_T,\nu_T,\mu_T)$ may be calibrated separately for each
maturity.  Each slice can then be internally arbitrage-consistent, but an
independently calibrated collection is not automatically free of calendar
arbitrage and is not the marginal family of a common dynamic process.
\end{remark}

\subsection{European calls and puts}

Let $K\ge0$ and set
\begin{equation}
  y_T(K):=\frac{K}{\bar S_0(T)}-A_T.
  \label{eq:exercise_boundary}
\end{equation}
Because $y_T(K)-x_T^*=K/\bar S_0(T)\ge0$, the call exercise region lies above
the default boundary.

\begin{proposition}[Closed-form valuation]
The time-zero call and put prices are
\begin{align}
  C_0(K,T)
  &=D_T\left[
      -\bar S_0(T)H_T(y_T)
      +(\bar S_0(T)A_T-K)(1-\cdf_T(y_T))
    \right],
  \label{eq:call_price}\\
  P_0(K,T)
  &=D_T\left[
      (K-\bar S_0(T)A_T)
      (\cdf_T(y_T)-\cdf_T(x_T^*))
    \right.
  \notag\\
  &\hspace{2.6cm}\left.
      {}-\bar S_0(T)(H_T(y_T)-H_T(x_T^*))
      +K\cdf_T(x_T^*)
    \right].
  \label{eq:put_price}
\end{align}
They satisfy market put--call parity,
\begin{equation}
  C_0(K,T)-P_0(K,T)=D_T(F_T-K).
  \label{eq:put_call_parity}
\end{equation}
\end{proposition}

\begin{proof}
For calls, $S_T>K$ is equivalent to $X_T>y_T$, so direct integration gives
\eqref{eq:call_price}.  For puts, the payoff equals $K$ on
$X_T\le x_T^*$ and equals
$K-\bar S_0(T)(A_T+X_T)$ on $x_T^*<X_T<y_T$; splitting the integral at
$x_T^*$ gives \eqref{eq:put_price}.  Subtracting and applying
\eqref{eq:martingale_A} gives \eqref{eq:put_call_parity}.
\end{proof}

The separate term $K\cdf_T(x_T^*)$ in \eqref{eq:put_price} is essential: it is
the put payoff on the atom at zero.  Omitting it prices an unfloored model that
allows negative terminal asset values.

\section{Greeks and strike derivatives}

The following derivatives hold with $(D_T,A_T,\sigma_T,\nu_T)$ fixed.  Thus
they are model derivatives without recalibration.  Write
$\bar S_0=\bar S_0(T)$, $y=y_T(K)$, and suppress the maturity subscript on
$f_T$, $\cdf_T$, and $H_T$.

\subsection{Prepaid-forward delta and gamma}

Differentiation with respect to the prepaid forward gives
\begin{align}
  \Delta_C^{\mathrm{pf}}
  :=\frac{\partial C_0}{\partial \bar S_0}
  &=D_T\left[A_T(1-\cdf_T(y))-H_T(y)\right],
  \label{eq:call_delta}\\
  \Delta_P^{\mathrm{pf}}
  :=\frac{\partial P_0}{\partial \bar S_0}
  &=\Delta_C^{\mathrm{pf}}-1,
  \label{eq:put_delta}\\
  \Gamma_C^{\mathrm{pf}}
  =\Gamma_P^{\mathrm{pf}}
  :=\frac{\partial^2 C_0}{\partial \bar S_0^2}
  &=D_T\frac{K^2}{\bar S_0^3}f_T(y).
  \label{eq:gamma}
\end{align}
The relation in \eqref{eq:put_delta} follows by differentiating
\eqref{eq:put_call_parity}, since $D_TF_T=\bar S_0$.

If deterministic dividend carry gives
$\bar S_0=S_0e^{-\delta T}$, then spot Greeks are
\begin{equation}
  \Delta_C^{S}=e^{-\delta T}\Delta_C^{\mathrm{pf}},\qquad
  \Gamma_C^{S}=e^{-2\delta T}\Gamma_C^{\mathrm{pf}},
  \label{eq:spot_greeks}
\end{equation}
and analogously for the put.

For the market forward supplied to the chain wrapper, $\bar S_0=D_TF_T$ and
$D_T$ is held fixed, so
\begin{equation}
  \Delta_C^{F}=D_T\Delta_C^{\mathrm{pf}},\qquad
  \Gamma_C^{F}=D_T^2\Gamma_C^{\mathrm{pf}},
  \label{eq:forward_greeks}
\end{equation}
and analogously for the put.

\subsection{Strike derivatives and digitals}

For $K>0$,
\begin{align}
  \frac{\partial C_0}{\partial K}
  &=-D_T[1-\cdf_T(y)],
  &
  \frac{\partial^2 C_0}{\partial K^2}
  &=D_T\frac{f_T(y)}{\bar S_0},
  \label{eq:call_strike_derivatives}\\
  \frac{\partial P_0}{\partial K}
  &=D_T\cdf_T(y),
  &
  \frac{\partial^2 P_0}{\partial K^2}
  &=D_T\frac{f_T(y)}{\bar S_0}.
  \label{eq:put_strike_derivatives}
\end{align}
Consequently the cash-or-nothing call is
\begin{equation}
  \mathrm{DigitalCall}_0(K,T)=D_T[1-\cdf_T(y)].
\end{equation}
The atom at $S_T=0$ creates a boundary effect at $K=0$; the derivatives above
are stated for strictly positive strikes.

\subsection{Volatility and tail-parameter sensitivities}

The martingale compensation $A_T$ changes when $\sigma_T$ or $\nu_T$ changes.
For a volatility bump, $X_T$ is proportional to $\sigma_T$.  Write
$x^*=-A_T$, $Q^*=1-\cdf_T(x^*)$, and $Q_y=1-\cdf_T(y)$.  Implicit
differentiation of \eqref{eq:martingale_A} gives
\begin{equation}
  \frac{\partial A_T}{\partial\sigma_T}
  =\frac{H_T(x^*)}{\sigma_TQ^*}.
  \label{eq:drift_vega_adjustment}
\end{equation}
The martingale-consistent call and put vegas are therefore identical:
\begin{equation}
  \boxed{
  \frac{\partial C_0}{\partial\sigma_T}
  =\frac{\partial P_0}{\partial\sigma_T}
  =\frac{D_T\bar S_0}{\sigma_T}
  \left[-H_T(y)+H_T(x^*)\frac{Q_y}{Q^*}\right].
  }
  \label{eq:martingale_consistent_vega}
\end{equation}
This is vega with respect to the prepaid-forward-normalized code parameter.
For the forward-normalized parameter in \eqref{eq:forward_normalization}, use
\begin{equation}
  \frac{\partial C_0}{\partial\widehat\sigma_T}
  =D_T^{-1}\frac{\partial C_0}{\partial\sigma_T},
  \qquad
  \frac{\partial P_0}{\partial\widehat\sigma_T}
  =D_T^{-1}\frac{\partial P_0}{\partial\sigma_T},
  \label{eq:forward_normalized_vega}
\end{equation}
with $D_T$ fixed.

More generally, let $\vartheta$ denote a distribution parameter, such as
$\nu_T$, and define
\begin{equation}
  G(A,\vartheta)
  :=\E_\vartheta[(A+X_\vartheta)^+]-D_T^{-1}=0.
\end{equation}
Implicit differentiation gives
\begin{equation}
  \frac{\partial A_T}{\partial\vartheta}
  =-\frac{\partial_\vartheta G(A_T,\vartheta)}
  {1-\cdf_{T,\vartheta}(-A_T)}.
  \label{eq:implicit_drift_greek}
\end{equation}
Therefore any numerical volatility or $\nu$ bump must re-solve
\eqref{eq:martingale_A}.  Holding $A_T$ fixed during the bump breaks the
forward identity.  The current implementation does not expose analytic Greeks;
the formulas above are adoption targets.  A consistent finite-difference Greek
must call \texttt{imply\_drift\_tdist} at every bumped parameter value.

\section{Stationary laws of stochastic-variance diffusions}

\subsection{General zero-current density}

For a positive one-dimensional diffusion
\begin{equation}
  \dd V_t=b(V_t)\dd t+a(V_t)\dd W_t,
\end{equation}
a zero-current invariant density, when it is normalizable, is
\begin{equation}
  \pi(v)\propto \frac{1}{a^2(v)}
  \exp\left(\int^v\frac{2b(u)}{a^2(u)}\dd u\right).
  \label{eq:stationary_general}
\end{equation}

\subsection{Multiplicative variance diffusion}

Consider
\begin{equation}
  \dd V_t=\kappa_m(\theta_m^2-V_t)\dd t
  +\varepsilon_mV_t\dd W_t.
  \label{eq:multiplicative_variance}
\end{equation}
Application of \eqref{eq:stationary_general} gives
\begin{equation}
  V_\infty\sim\IG(\alpha_m,\beta_m),\qquad
  \alpha_m=1+\frac{2\kappa_m}{\varepsilon_m^2},\qquad
  \beta_m=\frac{2\kappa_m\theta_m^2}{\varepsilon_m^2}.
  \label{eq:multiplicative_stationary}
\end{equation}
This is the inverse-gamma stationary law of the continuous-time GARCH
diffusion associated with Nelson's diffusion approximation
\cite{Nelson1990}.
Its stationary mean is $\theta_m^2$.  Its stationary variance is finite only
when $2\kappa_m>\varepsilon_m^2$, in which case
\begin{equation}
  \Var[V_\infty]
  =\frac{\theta_m^4\varepsilon_m^2}
  {2\kappa_m-\varepsilon_m^2}.
  \label{eq:multiplicative_stationary_variance}
\end{equation}

If the static mixture \eqref{eq:normal_ig_mixture} uses this stationary law,
then
\begin{equation}
  \nu=2\alpha_m=2+\frac{4\kappa_m}{\varepsilon_m^2},
  \qquad
  s^2=\frac{c\beta_m}{\alpha_m}.
  \label{eq:multiplicative_student_mapping}
\end{equation}
For $c=T$ and the code variance convention, $\sigma^2=\theta_m^2$.

\subsection{The \texorpdfstring{$3/2$}{3/2} variance diffusion}

Consider
\begin{equation}
  \dd V_t=\kappa_3V_t(\theta_3^2-V_t)\dd t
  +\varepsilon_3V_t^{3/2}\dd W_t.
  \label{eq:three_halves_variance}
\end{equation}
Its stationary law is
\begin{equation}
  V_\infty\sim\IG(\alpha_3,\beta_3),\qquad
  \alpha_3=2+\frac{2\kappa_3}{\varepsilon_3^2},\qquad
  \beta_3=\frac{2\kappa_3\theta_3^2}{\varepsilon_3^2}.
  \label{eq:three_halves_stationary}
\end{equation}
The reciprocal transformation maps this specification to a square-root
diffusion; the $3/2$ structure is represented in the nonlinear model of Ahn
and Gao \cite{AhnGao1999}.
Here $\theta_3^2$ is a drift level rather than the stationary mean:
\begin{align}
  \E[V_\infty]
  &=\frac{2\kappa_3\theta_3^2}{2\kappa_3+\varepsilon_3^2},
  \label{eq:three_halves_mean}
  \\
  \Var[V_\infty]
  &=\frac{2\kappa_3\varepsilon_3^2\theta_3^4}
  {(2\kappa_3+\varepsilon_3^2)^2}.
  \label{eq:three_halves_stationary_variance}
\end{align}
The corresponding static mixture is Student with $\nu=2\alpha_3$ and
$s^2=c\beta_3/\alpha_3$.  To match the code variance convention one must use
$\sigma^2=\E[V_\infty]$, not $\theta_3^2$.

\subsection{Heston/CIR variance diffusion}

Consider
\begin{equation}
  \dd V_t=\kappa_h(\theta_h^2-V_t)\dd t
  +\varepsilon_h\sqrt{V_t}\dd W_t.
  \label{eq:heston_variance}
\end{equation}
Its stationary law is gamma:
\begin{equation}
  V_\infty\sim\GammaDist(\alpha_h,\lambda_h),\qquad
  \alpha_h=\frac{2\kappa_h\theta_h^2}{\varepsilon_h^2},\qquad
  \lambda_h=\frac{2\kappa_h}{\varepsilon_h^2}.
  \label{eq:heston_stationary}
\end{equation}
This gamma invariant law is the stationary law of the CIR square-root process
\cite{CIR1985} used for variance in the Heston model \cite{Heston1993}.
The gamma scale is $\eta_h=1/\lambda_h=\varepsilon_h^2/(2\kappa_h)$ and
\begin{equation}
  \E[V_\infty]=\theta_h^2,
  \qquad
  \Var[V_\infty]
  =\frac{\theta_h^2\varepsilon_h^2}{2\kappa_h}.
  \label{eq:heston_stationary_moments}
\end{equation}
The stationary Laplace transform is
\begin{equation}
  L_V(\ell)=\left(\frac{\lambda_h}{\lambda_h+\ell}\right)^{\alpha_h}
  =(1+\eta_h\ell)^{-\alpha_h}.
  \label{eq:heston_laplace}
\end{equation}
The Feller condition $2\kappa_h\theta_h^2\ge\varepsilon_h^2$ makes zero
unattainable; the invariant gamma law may still exist under the standard CIR
boundary specification when the strict condition fails.

A static normal mixture with this gamma variance is not Student-$t$.  It has
the symmetric Bessel-$K$ density
\begin{equation}
  f_X(x)=
  \frac{2|x|^{\alpha_h-1/2}}
  {\sqrt{\pi}\,\Gamma(\alpha_h)(2c\eta_h)^{(\alpha_h+1/2)/2}}
  K_{\alpha_h-1/2}
  \left(\frac{2|x|}{\sqrt{2c\eta_h}}\right),
  \label{eq:gamma_normal_mixture}
\end{equation}
with characteristic function
\begin{equation}
  \phi_X(u)=\left(1+\frac{c\eta_hu^2}{2}\right)^{-\alpha_h}.
  \label{eq:gamma_mixture_cf}
\end{equation}
Unlike the Student mixture, this distribution has a moment-generating function
on a neighborhood of the origin:
\begin{equation}
  M_X(t)=\left(1-\frac{c\eta_ht^2}{2}\right)^{-\alpha_h},
  \qquad |t|<\sqrt{\frac{2}{c\eta_h}}.
  \label{eq:gamma_mixture_mgf}
\end{equation}
At the origin its density is finite if and only if $\alpha_h>1/2$, in which
case
\begin{equation}
  f_X(0)=\frac{\Gamma(\alpha_h-\tfrac12)}
  {\Gamma(\alpha_h)\sqrt{2\pi c\eta_h}}.
  \label{eq:gamma_mixture_origin}
\end{equation}
For $\alpha_h=1/2$ the singularity is logarithmic, and for $\alpha_h<1/2$ it
is a power singularity.

\subsection{Static mixture versus dynamic return law}

Suppose the asset diffusion is
\begin{equation}
  \frac{\dd S_t}{S_t}=\mu_t\dd t+\sqrt{V_t}\dd W_t^S.
\end{equation}
Conditionally on the variance path and with zero leverage correlation, the
centered stochastic integral satisfies
\begin{equation}
  \int_0^T\sqrt{V_t}\dd W_t^S
  \ \Big|\ \{V_t:0\le t\le T\}
  \sim N\left(0,\int_0^TV_t\dd t\right).
  \label{eq:integrated_variance_mixture}
\end{equation}
Thus the exact finite-horizon mixing law uses integrated variance rather than
the stationary marginal $V_\infty$.  In addition,
\begin{equation}
  \dd\log S_t=\left(\mu_t-\frac12V_t\right)\dd t
  +\sqrt{V_t}\dd W_t^S,
\end{equation}
so log returns have a variance-dependent conditional mean.  The stationary
mixtures above should therefore be described as exact one-step hierarchical
models, short-horizon frozen-variance approximations, or independent terminal
specifications.  They are not, without further work, exact terminal laws of the
continuous-time diffusions.

The same distinction appears in dynamic option-pricing work.  For example,
Barone-Adesi, Rasmussen, and Ravanelli
\cite{BaroneAdesiRasmussenRavanelli2005} derive moments of integrated variance
for the GARCH diffusion and use them in an analytical approximation to
European option values; they do not replace integrated variance by an
independent draw from its stationary marginal.

\section{Correspondence with the implementation}

The existing library pricer implements the symmetric $q=0$ formulation.
The chapter script \texttt{risk\_premium\_smile\_examples.py} implements the
new $p,\eta,q$ extension through conditional-normal quadrature and generates
Figures~\ref{fig:risk_premium_smiles}--\ref{fig:p_fixed_variance}.  The
following table records the symmetric formulation adopted by the existing
library code.
\begin{center}
\begin{tabular}{@{}ll@{}}
\toprule
Mathematical object & Implementation object \\
\midrule
$T$ & \texttt{ttm} \\
$D_T$ & \texttt{discfactor} or $\exp(-\texttt{rf\_rate}\,T)$ \\
$F_T$ & \texttt{forward} in the chain wrapper \\
$\bar S_0=D_TF_T$ & argument \texttt{spot} in the chain wrapper \\
$\sigma$ (prepaid-forward normalization) & \texttt{vol} \\
$\nu$ & \texttt{nu} \\
$s_T$ in \eqref{eq:code_student_scale} & \texttt{compute\_upsilon} \\
$\cdf_T(x)$ & \texttt{cdf\_tdist(x, mu=0, ...)} \\
$H_T(x)$ & \texttt{cum\_mean\_tdist(x, mu=0, ...)} \\
$\mu_T$ & \texttt{risk\_neutral\_mu} or \texttt{drift} \\
$A_T=1+\mu_TT$ & \texttt{1.0 + risk\_neutral\_mu*ttm} \\
\bottomrule
\end{tabular}
\end{center}

For a one-maturity slice whose stored drift was solved at the same maturity,
the implementation contract is:
\begin{enumerate}
  \item set $\bar S_0=D_TF_T$ and
  $r_T=-\log(D_T)/T$;
  \item compute $s_T$ from \eqref{eq:code_student_scale};
  \item solve \eqref{eq:martingale_mu} for $\mu_T$;
  \item evaluate \eqref{eq:call_price} or \eqref{eq:put_price};
  \item verify \eqref{eq:put_call_parity}.
\end{enumerate}

The chain wrapper follows this contract exactly for a one-maturity slice with
matching \texttt{params.ttm} and \texttt{params.drift}.  Although
\texttt{price\_chain} accepts arrays of maturities, it passes one stored drift
to every maturity and does not enforce that match.  The standalone function
\texttt{compute\_vanilla\_price\_tdist} may interpret its argument named
\texttt{spot} as the actual cash spot only when carry is absent or separately
embedded; the chain wrapper removes the ambiguity by passing $D_TF_T$.

The valuation core already implements the centered call and put formulas in
this note.  Full code adoption nevertheless requires the following corrections
and guardrails.
\begin{enumerate}
  \item The model shock is an arithmetic relative-return shock, not a log
  return or a Black--Scholes volatility.  Student exponential moments do not
  exist, which is precisely why \eqref{eq:terminal_asset} is used instead of an
  exponential price map.  The existing ``terminal log-returns'' descriptions
  should be corrected, as should the ``normal implied vol'' description of the
  Student-$t$ inversion routine.
  \item The public density helpers accept a non-zero location parameter.  The
  terminal pricer always calls the CDF and partial-moment helpers with location
  zero and applies $\mu_T$ externally.  That centered route is exact.  For a
  non-zero location $m=\mu T$, however, the correct partial moment is
  \[
    H_Y(x)=m\cdf_Y(x)+sH_Z((x-m)/s).
  \]
  The present \texttt{cum\_mean\_tdist} uses \texttt{mu*cdf} instead of
  \texttt{mu*ttm*cdf}; this latent helper bug does not affect centered pricing.
  \item The quantity $\mu_T=(A_T-1)/T$ is a maturity-specific additive
  location compensation, not the drift coefficient of a continuous-time
  return process.  A different maturity generally requires a different root.
  \item Calling \texttt{compute\_vanilla\_price\_tdist} with
  \texttt{is\_compute\_risk\_neutral\_mu=False} and no explicit drift sets
  $\mu_T$ equal to the risk-free rate.  Because the floor itself has value,
  that fallback does not satisfy \eqref{eq:martingale_A} and should be rejected
  or clearly marked non-risk-neutral.
  \item The raw pricer currently treats option types \texttt{IC} and
  \texttt{IP} as ordinary calls and puts.  Elsewhere those labels denote
  inverse payoffs.  The inverse put genuinely diverges at the
  positive-probability atom $S_T=0$.  For $K>0$, the inverse call has a
  removable $0/0$ value and continuous extension zero there, but it still
  requires an explicit zero-state convention and an implementation guard.
  The model should accept only \texttt{C}/\texttt{P} until those contracts are
  defined separately.
  \item Calibration is performed one maturity at a time.  A stationary SDE
  interpretation would impose cross-maturity parameter restrictions that the
  current calibrator does not impose.
\end{enumerate}

\section{Verification identities}

Any implementation of this formulation should verify the following independently.
\begin{enumerate}
  \item Distribution normalization and variance:
  \[
    \int_{-\infty}^{\infty}f_T(x)\dd x=1,
    \qquad
    \int_{-\infty}^{\infty}x^2f_T(x)\dd x=\sigma^2T.
  \]
  \item Martingale and default atom:
  \[
    \bar S_0\{A_T[1-\cdf_T(-A_T)]-H_T(-A_T)\}=F_T,
    \qquad p_T^{(0)}=\cdf_T(-A_T).
  \]
  \item Put--call parity \eqref{eq:put_call_parity} for several strikes.
  \item Finite-difference agreement with \eqref{eq:call_delta}--\eqref{eq:gamma}
  while holding calibrated parameters fixed.
  \item Agreement of the analytic Student density with an independent
  inverse-gamma/normal mixture simulation.
  \item Normalization and the weighted martingale identities for the
  physical-to-pricing kernel:
  \[
    \E_{\mathbb P}[\Lambda_{p,\eta,q}]=1,\qquad
    \E_{\mathbb P}[\Lambda_{p,\eta,q}S_T]=F_T.
  \]
  \item At $q=0$, agreement of
  \eqref{eq:skew_call_price} with the independent Student partial-moment
  formula; at $q\ne0$, agreement with adaptive integration of the gamma
  precision density and with direct integration of the marginal Bessel-$K$
  density \eqref{eq:skew_t_density}.
  \item Put--call parity, Black implied-volatility round trips, decreasing
  call prices, non-negative strike convexity, and convergence under
  quadrature refinement.
  \item For \eqref{eq:variance_neutral_eta}, preservation of
  $\E_{\mathbb Q}[V]$, and equality of all duplicate baseline curves.
\end{enumerate}

The accompanying deterministic study uses 256 generalized Laguerre nodes.
Across every plotted strike and scenario, the maximum absolute errors are
\begin{center}
\begin{tabular}{@{}lr@{}}
\toprule
Check & Maximum absolute error \\
\midrule
Martingale solver residual & $3.06\times10^{-14}$ \\
Independent martingale check & $2.80\times10^{-8}$ \\
Put--call parity & $3.04\times10^{-14}$ \\
Symmetric Student closed form & $2.26\times10^{-8}$ \\
Adaptive gamma-precision integration & $2.80\times10^{-8}$ \\
Marginal Bessel-$K$ integration & $2.80\times10^{-8}$ \\
Black price--volatility round trip & $3.33\times10^{-16}$ \\
128 versus 256 quadrature nodes & $1.75\times10^{-7}$ \\
Call monotonicity violation & $0$ \\
Call convexity violation & $0$ \\
Variance-neutral mean-$V$ error & $0$ \\
Duplicate-baseline curve error & $0$ \\
\bottomrule
\end{tabular}
\end{center}
The comparisons use different numerical representations, rather than merely
re-evaluating the same quadrature formula.

\section{Conclusion}

The inverse-gamma variance mixture provides a precise static foundation for a
Student-$t$ return shock.  European option valuation becomes arbitrage-consistent
at a fixed maturity after three additions: use the prepaid forward as the price
scale, floor the terminal asset at zero, and solve the drift from the martingale
condition.  The resulting exact put formula includes the atom at zero, and the
Greeks depend on both exercise probability and truncated first moments.

The recommended physical-to-pricing change of measure is
\eqref{eq:combined_pricing_kernel}.  It avoids the invalid marginal Student
Esscher transform, preserves the inverse-gamma family through $(p,\eta)$,
and generates asymmetric return risk through the conditionally normalized
parameter $q$.  In smile coordinates, $\eta$ primarily controls variance
level, variance-neutral $p$ controls tail curvature, and $q$ controls skew.
Valuation reduces to one-dimensional conditional-Bachelier integration, and
the risk-premium Greeks follow from centered inverse-gamma scores plus the
implicit martingale-location response.

Multiplicative and $3/2$ variance diffusions have inverse-gamma stationary laws
and therefore motivate static Student mixtures.  Heston variance has a gamma
stationary law and motivates a Bessel-$K$, not Student, mixture.  These
stationary constructions should remain conceptually separate from exact
finite-horizon continuous-time return laws, which depend on integrated variance.
The terminal kernel is a valid one-period pricing law; a dynamic interpretation
requires a true Girsanov density, boundary admissibility, and an explicit
criterion selecting one equivalent martingale measure from the incomplete
market.

\begin{thebibliography}{99}
\addcontentsline{toc}{section}{References}
\small
\raggedright

\bibitem{Praetz1972}
P.~D. Praetz,
``The distribution of share price changes,''
\emph{The Journal of Business}, vol.~45, no.~1, pp.~49--55, 1972.
\doi{10.1086/295425}.

\bibitem{Pinn2000}
K. Pinn,
``Minimal variance hedging of options with Student-$t$ underlying,''
\emph{Physica A: Statistical Mechanics and its Applications},
vol.~276, nos.~3--4, pp.~581--595, 2000.
\doi{10.1016/S0378-4371(99)00498-7}.

\bibitem{Vince1992}
R. Vince,
\emph{The Mathematics of Money Management: Risk Analysis Techniques for
Traders}. New York: John Wiley \& Sons, 1992. ISBN 978-0-471-54738-9.

\bibitem{CassidyHampOuyed2010}
D.~T. Cassidy, M.~J. Hamp, and R. Ouyed,
``Pricing European options with a log Student's $t$-distribution: A Gosset
formula,''
\emph{Physica A: Statistical Mechanics and its Applications},
vol.~389, no.~24, pp.~5736--5748, 2010.
\doi{10.1016/j.physa.2010.08.037}.

\bibitem{CassidyHampOuyed2013}
D.~T. Cassidy, M.~J. Hamp, and R. Ouyed,
``Log Student's $t$-distribution-based option sensitivities: Greeks for the
Gosset formulae,''
\emph{Quantitative Finance}, vol.~13, no.~8, pp.~1289--1302, 2013.
\doi{10.1080/14697688.2012.744087}.

\bibitem{CastelliLeonenkoShchestyuk2017}
F. Castelli, N.~N. Leonenko, and N. Shchestyuk,
``Student-like models for risky asset with dependence,''
\emph{Stochastic Analysis and Applications},
vol.~35, no.~3, pp.~452--464, 2017.
\doi{10.1080/07362994.2016.1266945}.

\bibitem{LeonenkoLiuShchestyuk2025}
N.~N. Leonenko, A. Liu, and N. Shchestyuk,
``Student models for a risky asset with dependence: Option pricing and
Greeks,''
\emph{Austrian Journal of Statistics},
vol.~54, no.~1, pp.~138--165, 2025.
\doi{10.17713/ajs.v54i1.1952}.

\bibitem{Ivanov2023}
R.~V. Ivanov,
``On the stochastic volatility in the generalized Black--Scholes--Merton
model,''
\emph{Risks}, vol.~11, no.~6, article~111, 2023.
\doi{10.3390/risks11060111}.

\bibitem{AasHaff2006}
K. Aas and I. Hob\ae k Haff,
\textquotedblleft The generalized hyperbolic skew Student's
$t$-distribution,\textquotedblright
\emph{Journal of Financial Econometrics},
vol.~4, no.~2, pp.~275--309, 2006.
\doi{10.1093/jjfinec/nbj006}.

\bibitem{ChoulliStrickerLi2007}
T. Choulli, C. Stricker, and J. Li,
\textquotedblleft Minimal Hellinger martingale measures of order
$q$,\textquotedblright
\emph{Finance and Stochastics}, vol.~11, no.~3, pp.~399--427, 2007.
\doi{10.1007/s00780-007-0039-3}.

\bibitem{JeanblancKloppelMiyahara2007}
M. Jeanblanc, S. Kl\"oppel, and Y. Miyahara,
\textquotedblleft Minimal $f^q$-martingale measures for exponential
L\'evy processes,\textquotedblright
\emph{The Annals of Applied Probability},
vol.~17, nos.~5--6, pp.~1615--1638, 2007.
\doi{10.1214/07-AAP439}.

\bibitem{DesmettreLeobacherRogers2021}
S. Desmettre, G. Leobacher, and L.~C.~G. Rogers,
\textquotedblleft Change of drift in one-dimensional
diffusions,\textquotedblright
\emph{Finance and Stochastics}, vol.~25, no.~2, pp.~359--381, 2021.
\doi{10.1007/s00780-021-00451-w}.

\bibitem{Nelson1990}
D.~B. Nelson,
``ARCH models as diffusion approximations,''
\emph{Journal of Econometrics}, vol.~45, nos.~1--2, pp.~7--38, 1990.
\doi{10.1016/0304-4076(90)90092-8}.

\bibitem{AhnGao1999}
D.-H. Ahn and B. Gao,
``A parametric nonlinear model of term structure dynamics,''
\emph{The Review of Financial Studies},
vol.~12, no.~4, pp.~721--762, 1999.
\doi{10.1093/rfs/12.4.721}.

\bibitem{CIR1985}
J.~C. Cox, J.~E. Ingersoll, Jr., and S.~A. Ross,
``A theory of the term structure of interest rates,''
\emph{Econometrica}, vol.~53, no.~2, pp.~385--407, 1985.
\doi{10.2307/1911242}.

\bibitem{Heston1993}
S.~L. Heston,
``A closed-form solution for options with stochastic volatility with
applications to bond and currency options,''
\emph{The Review of Financial Studies},
vol.~6, no.~2, pp.~327--343, 1993.
\doi{10.1093/rfs/6.2.327}.

\bibitem{BaroneAdesiRasmussenRavanelli2005}
G. Barone-Adesi, H. Rasmussen, and C. Ravanelli,
``An option pricing formula for the GARCH diffusion model,''
\emph{Computational Statistics \& Data Analysis},
vol.~49, no.~2, pp.~287--310, 2005.
\doi{10.1016/j.csda.2004.05.014}.

\end{thebibliography}

\end{document}
